\begin{pastexam}[label={px:ncku-115-prog-14}]{成大資訊聯招 115 學年度 演算法問答第 1 題}
（問答）

[10\%] Assuming that you are a medical information analyst, currently analyzing an
algorithm designed to process a 3D MRI volume. The input is a cubic volume of data with
side length n (total voxels = n $\times$ n $\times$ n). Assume n is a power of 2.

Pseudo-code:
\begin{center}
\fbox{\begin{minipage}{0.8\linewidth}
\textbf{Surface\_Scan(n)}\\
\hspace*{1.5em}total\_intensity = 0\\
\hspace*{1.5em}for i = 1 to n\\
\hspace*{3em}for j = 1 to 100\\
\hspace*{4.5em}total\_intensity = total\_intensity + Read\_Pixel(i, j)\\
\hspace*{1.5em}return total\_intensity

\medskip
\textbf{MRI\_Voxel\_Process(n)}\\
\hspace*{1.5em}if n <= 1\\
\hspace*{3em}return

\smallskip
\hspace*{1.5em}for k = 1 to n\\
\hspace*{3em}Surface\_Scan(n)

\smallskip
\hspace*{1.5em}for x = 0 to 1\\
\hspace*{3em}for y = 0 to 1\\
\hspace*{4.5em}for z = 0 to 1\\
\hspace*{6em}MRI\_Voxel\_Process(n/2)
\end{minipage}}
\end{center}

Note: Assume that the function Read\_Pixel(i, j) takes constant time

\underline{\textbf{Questions}}:
\begin{enumerate}[label=(\arabic*)]
\item {[3\%]} Analyze the \textbf{Surface\_Scan(n)} function and determine its time
complexity using $\Theta$-notation.
\begin{itemize}\item Time complexity: \underline{\hspace{5cm}}\end{itemize}
\item {[3\%]} Write down the recurrence relation for the running time T(n) of
\textbf{MRI\_Voxel\_Process(n)}.
\begin{itemize}\item T(n) = \underline{\hspace{5cm}}\end{itemize}
\item {[4\%]} Analyze the asymptotic growth rate of the algorithm, and determine the
overall asymptotic tight bound ($\Theta$) for the running time of
\textbf{MRI\_Voxel\_Process(n)}. Please write down your calculation process.
\begin{itemize}
\item Calculation process: \underline{\hspace{5cm}}
\item Overall asymptotic tight bound: \underline{\hspace{4cm}}
\end{itemize}
\end{enumerate}
\end{pastexam}

\begin{pastsoln}
答案：(1) $\Theta(n)$；(2) $T(n) = 8T(n/2) + \Theta(n^2)$，$T(1) = \Theta(1)$；(3) $\Theta(n^3)$。

(1) $\Theta(n)$。 \textsf{Surface\_Scan} 的外層跑 $n$ 次，內層固定跑 $100$ 次，
\textsf{Read\_Pixel} 一共呼叫 $100n$ 次，每次常數時間，再加上常數個其他步驟。 內層的
上界是常數 $100$，不是 $n$，所以不是 $n^2$。

(2) 照資結書第~2.1~節的四件事來寫。 規模是 $n$；基礎情況 $n \le 1$ 直接返
回，$\Theta(1)$；遞迴呼叫：$x$、$y$、$z$ 各有 $0$、$1$ 兩個值，共 $2 \cdot 2 \cdot 2
= 8$ 次，每次規模 $n/2$；呼叫以外的工作：呼叫 \textsf{Surface\_Scan}$(n)$ 共 $n$ 次，
每次 $\Theta(n)$，合計 $\Theta(n^2)$。 所以
\[
T(n) = 8T(n/2) + \Theta(n^2) \quad (n \ge 2), \qquad T(1) = \Theta(1).
\]

(3) $\Theta(n^3)$。 用主定理（資結書定理~2.4）：$a = 8$、$b = 2$，
$n^{\log_2 8} = n^3$；$f(n) = \Theta(n^2) = O(n^{3 - \varepsilon})$，取
$\varepsilon = 1$，是情況 1，$T(n) = \Theta(n^3)$。

也可以把常數算出來。 只數 \textsf{Read\_Pixel} 的呼叫次數 $R(n)$：$R(1) = 0$，
$R(n) = 8R(n/2) + 100n^2$。 用資結書引理~2.6 展開，第 $j$ 層是
$8^j \cdot 100 (n/2^j)^2 = 100 n^2 \cdot 2^j$，
\[
R(n) = 100 n^2 \sum_{j=0}^{\lg n - 1} 2^j = 100 n^2 (n - 1) = 100 n^3 - 100 n^2 .
\]
每層是上一層的 $2$ 倍，越來越大，最下面那層決定答案（資結書第~2.3~節的
「葉子贏」）。 直覺上也合理：整個體積有 $n^3$ 個 voxel，八個子問題的 voxel 加起來還
是 $n^3$；每一層的表面掃描只占它那一層的一小部分。
\end{pastsoln}
